%!TEX root = mainQlin.tex


\section{Mechanisms}

In this section, we overview several of the key ideas behind the regret bound in Theorem \ref{thm:main}. We defer the full proof of Theorem \ref{thm:main} and Theorem \ref{thm:main_mis} to Appendix \ref{app:proof_main} and Appendix \ref{app:proof_main_mis} respectively.

In Section \ref{sec:results}, we mentioned that the LSVI algorithm is motivated from the Bellman optimality equation Eq.~\eqref{eq:opt_bellman}. It remains to verify that line \ref{line:w} in Algorithm \ref{algo:LSVI-UCB} indeed well approximates the Bellman optimality equation, 
which turns out to require not only the linear MDP structure but also hinges on several other facts. 

To simplify our presentation, in this section we treat the regularization parameter $\lambda$ loosely as being sufficiently small so that $\Lambda_h^{-1} \sum_{\tau =1}^{k-1}  \bphi(x^{\tau}_h,  a^{\tau}_h)\bphi(x^{\tau}_h, a^{\tau}_h)\trans \approx \I$. We will focus in this section on a fixed episode $k$, and drop the dependency of parameters and value functions on $k$ when it is clear from the context. Now, ignoring the UCB bonus, the least-squares solution (line \ref{line:w}) gives the following estimate of the action-value function:
\begin{equation*}
Q_h(x, a) \approx \bphi(x, a)\trans \w_h = \bphi(x, a)\trans\Lambda_h^{-1} \sum_{\tau=1}^{k-1} \bphi(x^{\tau}_h,  a^{\tau}_h) [r_h(x^{\tau}_h, a^{\tau}_h) +  V_{h+1}(x^{\tau}_{h+1})],
\end{equation*}
where $V_{h+1}(\cdot) = \max_{a\in\cA} Q_{h+1}(\cdot, a)$. Plugging in $r_h(\cdot, \cdot) = \bphi(\cdot, \cdot)\trans \btheta_h$, we know the first term on the right-hand side approximates $r_h(x, a)$. Comparing this to Eq.~\eqref{eq:opt_bellman}, it remains to show why the second term of right-hand side approximates $\P_h V_{h+1}(x, a)$. We thus define our empirical Markov transition measure as
\begin{equation*}
\hat{\P}_h (\cdot | x, a) \defeq \bphi(x, a)\trans\Lambda_h^{-1} \sum_{\tau=1}^{k-1} \bphi(x^{\tau}_h,  a^{\tau}_h) \delta(\cdot, x^{\tau}_{h+1}),
\end{equation*}
where the $\delta$-measure $\delta(\cdot, x)$ puts an atom on element $x$. It remains to verify that $\hat{\P}_h V_{h+1}(x, a) \approx \P_h V_{h+1}(x, a)$.
% to better explain the intuition behind our theoretical results, we introduce the core mechanisms that enable us to establish the regret upper bound in Theorem \ref{thm:main}. 
To establish this, we use a measure $\bar{\P}_h$ to bridge these two quantities: 
\begin{equation} \label{eq:P_bar}
\bar{\P}_h (\cdot | x, a) \defeq \bphi(x, a)\trans\Lambda_h^{-1} \sum_{\tau=1}^{k-1} \bphi(x^{\tau}_h,  a^{\tau}_h) \P_h(\cdot | x^{\tau}_h,  a^{\tau}_h).
\end{equation}
Our analysis depends on the following two key steps.

\paragraph{Step 1: Prove $\hat{\P}_h V_{h+1}(x, a) \approx \bar{\P}_h V_{h+1}(x, a)$ via  Value-Aware Uniform Concentration.} Computing the difference, we have $(\hat{\P}_h - \bar{\P}_h)V_{h+1}(x, a) 
= \bphi(x, a)\trans\Lambda_h^{-1} \sum_{\tau=1}^{k-1} \bphi(x^{\tau}_h,  a^{\tau}_h) [V_{h+1}(x^\tau_{h+1}) - \P_h V_{h+1}(x^{\tau}_h,  a^{\tau}_h)]$. Since $\x^\tau_{h+1}$ is a sample from the distribution $\P_h(\cdot|x^{\tau}_h,  a^{\tau}_h)$, we would expect this term to be small due to concentration. This would be the case if function $V_{h+1}$ is \emph{fixed and independent of the samples} $\{x^\tau_{h+1}\}_{\tau=1}^{k-1}$. Then, $V_{h+1}(x^\tau_{h+1}) - \P_h V_{h+1}(x^{\tau}_h,  a^{\tau}_h)$ is a zero-mean random variable in $[-H, H]$, and we could aim to use a concentration inequality for self-normalized processes to bound $(\hat{\P}_h - \bar{\P}_h)V_{h+1}(x, a)$. Please see Theorem \ref{thm:self_norm} or \cite{abbasi2011improved} for more detail on this approach.

However, the function $V_{h+1}$ in Algorithm \ref{algo:LSVI-UCB} is again computed by least-squares value iteration in later steps $[h+1, H]$ and it thus inevitably depends on the choices of actions $\{a^\tau_{h+1}\}_{\tau=1}^{k-1}$, and thus also samples $\{x^\tau_{h+1}\}_{\tau=1}^{k-1}$. Therefore, the concentration of self-normalized process does not apply directly. To resolve this issue, we establish the uniform concentration over all value functions in the following class:
\begin{equation}\label{eq:class_value}
\mathcal{V} = \Big \{ V(\cdot ) | V(\cdot ) = \min \bigl  \{ \max_{a \in \cA} \bphi (\cdot, a)  ^\top  \w + \beta \sqrt{\bphi (\cdot, a) \Lambda^{-1} \bphi (\cdot, a)}, H\bigr \}, \w \in \R^d, \beta \in \R, \Lambda \in \R^{d\times d} \Bigr \},
\end{equation}
where the parameters $\w, \beta, \Lambda$ are all bounded. We ensure that Algorithm \ref{algo:LSVI-UCB} only uses value functions within this class $\mathcal{V}$, which has a reasonably small covering number. This gives, with high probability, $|(\hat{\P}_h - \bar{\P}_h)V_{h+1}(x, a)| \le \tilde{\cO}(dH)\cdot(\bphi(x, a) \Lambda^{-1}_h \bphi(x, a))^{1/2}$  (Lemma \ref{lem:stochastic_term}). 


\paragraph{Step 2: Show $\bar{\P}_h V_{h+1}(x, a) \approx \P_h V_{h+1}(x, a)$ due to Linear Markov Transitions.} One big challenge in RL with function approximation is that, due to the large state space, the learner may never visit the neighborhood of a state-action pair twice. This raises a question of how to use the experiences from other state-action pairs to infer information about a state-action pair of interest. In Eq.~\eqref{eq:P_bar}, $\bar{P}_h(\cdot|x, a)$ provides such an estimate via regularized least-squares. Our modeling assumption of a linear MDP (Assumption \ref{assumption:linear}) ensures that this least-square estimate is valid: since $\P_h(\cdot|x, a)=\bphi(x, a)\trans\bmu_h(\cdot)$ for any $(x, a)$ pair, we have
\begin{equation*}
\bar{\P}_h (\cdot | x, a) = \bphi(x, a)\trans\Lambda_h^{-1} \sum_{\tau=1}^{k-1} \bphi(x^{\tau}_h,  a^{\tau}_h)\bphi(x^{\tau}_h,  a^{\tau}_h)\trans\bmu_h(\cdot)
\approx \bphi(x, a)\trans\bmu_h(\cdot) = \P_h (\cdot | x, a).
\end{equation*}

In summary, combining step 1 and step 2, we establish $\hat{\P}_h V_{h+1}(x, a) \approx \P_h V_{h+1}(x, a)$, and hence show that LSVI approximates the optimal Bellman equation. We emphasize that despite being linear, the Markov transition model $\P_h(\cdot|x, a) = \bphi(x, a)\trans \bmu_h(\cdot)$ can still have infinite degrees of freedom since the measure $\bmu_h$ is unknown. Therefore, within a finite number of samples, no algorithm can establish that $\hat{\P}_h$ and $\P_h$ are close in total variation distance. In contrast, our algorithm only requires $\hat{\P}_h V_{h+1}(x, a) \approx \P_h V_{h+1}(x, a)$ for all value functions $V_{h+1}$ in a small function class $\mathcal{V}$ (especially in step 1). This bypasses the need for fully learning the transition model $\P_h$. Thus, our algorithm can also be viewed as ``\emph{model-free}'' in this sense.  

Finally, with the above key observations in mind, our proof proceeds by leveraging and adapting techniques from the literature on tabular MDP and linear bandits. Please see Appendix \ref{app:proof_main} and \ref{app:proof_main_mis} for the details.


% This is the place where the assumpton of linear MDP plays a key role. It allows the least square solution provides a good approximate to the expectation under true transition model.


% Note the most difference... (instead of bound TV, we bound the expected on the restrict class of functions.)

% Remaining, adapt tabular, linear bandits results.




% {\noindent \bf Mechanism 1: Linear Transition Model.}


%  The first mechanism is that the linear transition model enables us to perform approximate value iteration via least-squares regression. Specifically, for any $h \in [H]$, let $Q_{h+1}$ be the current estimate of the $h+1$-step action-value function. Then value iteration aims to update $Q_h $ by 
%  \#\label{eq:lin_bellman}
%  Q_h(x,a)  =  r(x, a) + (\P_h V _{h+1}  )(x , a) = \la \bphi(x, a), \btheta_h \ra + \bigg \la \bphi(x, a) , \int_{\cS} V_{h+1} (x')   \ud \bmu_h(x') \bigg\ra  ,
%  \#
% where $V_{h+1} (x ) = \max _{a\in \cA} Q_{h+1} (x,a)$ for all $x\in \cS$.
% Thus, when $K$ is sufficiently large, solving a least-squares regression problem using data $\{ r(x^{ k}_h, a^{k }_h) + V(x^{k}_{h+1}),\bphi( x^{ k}_h, a_h^{k } ) \}_{ k \in [K]} $ yields an accurate estimate of $Q_h$.
% Moreover, by \eqref{eq:lin_bellman}, we can construct an empirical estimator of the Markov transition operator by regressing $\{ V(x^{k}_{h+1}) \}_{k\in [K]} $ to  $\{\bphi( x^{ k}_h, a_h^{k } ) \}_{ k \in [K]} $.  When combined with $\ell_2$-regularization, the regression coefficient is $  \Lambda_h^{-1} \sum_{\tau=1}^{K} \bphi(x^{\tau}_h,  a^{\tau}_h) \cdot V_{h+1}(x^{\tau}_{h+1} )$, where $  \Lambda_h = \sum_{\tau =1}^{K}  \bphi(x^{\tau}_h,  a^{\tau}_h)\bphi(x^{\tau}_h, a^{\tau}_h)\trans + \lambda \cdot  \I $.  Hence, for any $V_{h+1}$, 
% \#\label{eq:empirical_bellman}
% ( \hat \P_h V _{h+1} ) (x,a) =  \biggl \la \bphi (x, a) , \Lambda_h^{-1} \sum_{\tau=1}^{K} \bphi(x^{\tau}_h,  a^{\tau}_h) \cdot V_{h+1}(x^{\tau}_{h+1} ) \biggr \ra  
% \#
% is an empirical estimator of $(\P_h V _{h+1})(x, a)$, which converges to   $\P_h V _{h+1}(x, a)$ when $K$ goes to infinity.
 
 
%  \iffalse 
% Talk about our empirical transition matrix estimator

% , linear transition -> enough samples -> true estimator 
% \fi 


% {\noindent \bf Mechanism 2: Value-Aware Concentration.}
% Our  second pivotal technique   is to bound the difference between $\hat \P_h V _{h+1} $  and $\P_h V _{h+1}$ via uniform concentration with respect to the family to all possible value functions.
% Note that the family of  value functions is a subclass   
% \$
% \cF = \Big \{ V(\cdot ) \colon V(\cdot ) = \min \bigl  \{ \max_{a \in \cA} \bphi (\cdot, a)  ^\top  \w, H   \bigr \}, \w \in \R^d \Bigr \} .\$
%  This approach is in stark contrast to directly establishing uniform concentration over function class $\cF$. Since parameter $\w \in \R^d$ in $\cF$ might not be bounded, function class  $\cF$ can have huge capacity. To overcome such a difficulty, we incorporate   $\ell_2$-regularization in least-squares value iteration  (line \ref{line:w} in Algorithm \ref{algo:LSVI-UCB}). 
%  As we will show in the appendix, with $\ell_2$-regularization, it suffices to only focus on value functions in  $\cF$ satisfying $\| \w \| \leq 2 H  \sqrt{d k / \lambda}$, where $\lambda $ is the regularization parameter.  For properly chosen $\lambda$, this family of value functions forms a much smaller subclass  of $\cF$ and thus enable us to establish sharper concentration results.

% \iffalse 
% Mechanism 2, Concentration, we only require empirical transition and population are similar for all value functions used in the algorithm, the complexity of which is under controls.
% \fi
